% TOP_PROBLEM: {"id": 30006809, "title": "Hopf conjecture for nonnegative sectional curvature", "category_id": 6, "category": {"id": 6, "name": "geometry", "display_name": "Geometry", "description": "Euclidean and non-Euclidean geometry, geometric structures.", "slug": "geometry", "order_index": 6, "created_at": "2026-07-31T15:26:25.671Z"}, "status": "open"}
% FIRST_POSTED: 2026-10-01
% ATTEMPT_STATUS: unresolved
% !TeX program = lualatex
% Definition and sources only; this file is not an LLM solution attempt.
\documentclass[11pt]{article}
\usepackage[margin=25mm]{geometry}
\usepackage{fontspec}
\setmainfont{Latin Modern Roman}
\usepackage{amsmath,amssymb,mathtools}
\usepackage{unicode-math}
\setmathfont{Latin Modern Math}
\usepackage{xurl,hyperref}
\hypersetup{colorlinks=true,urlcolor=blue}
\setcounter{secnumdepth}{0}
\setlength{\parindent}{0pt}
\setlength{\parskip}{5pt}
\setlength{\emergencystretch}{3em}
\begin{document}
\section*{\texorpdfstring{Hopf conjecture for nonnegative sectional curvature}{Hopf conjecture for nonnegative sectional curvature}}\label{30006809}
\subsection{Definitions and mathematical statement}
For a closed even-dimensional Riemannian manifold $M$, let $K_g(\sigma)$ be sectional curvature on a tangent two-plane, with the sign convention positive on round spheres, and let
$\chi(M)=\sum_i(-1)^i\dim H_i(M,{\mathbb{Q}})$. The assertion is
\[
 K_g(\sigma)\ge0\text{ on every tangent two-plane}
 \quad\Longrightarrow\quad\chi(M)\ge0.
\]
Zero sectional curvature and zero Euler characteristic are allowed. The manifold need not be simply connected or have an assumed isometric group action. This nonnegative-curvature statement is distinct from the strict positive-curvature version requiring strictly positive Euler characteristic, and neither is merely a scalar-curvature condition.
\par\smallskip\textit{Formulation and concept references:} \href{https://link.springer.com/article/10.1365/s13291-020-00225-x}{[S1]}.

\subsection{Short English statement}
Must every closed even-dimensional manifold with nonnegative sectional curvature have nonnegative Euler characteristic?
\subsection{Research attempt: surfaces, product closure, and a scalar-curvature stress test}
\textbf{Status.} We prove the selected nonnegative-curvature assertion in dimension two and for classes built by products and finite free isometric quotients from known examples. We also exhibit why replacing sectional curvature by scalar curvature invalidates the claim. The primary discussion [S1] supplies context for the higher-dimensional Hopf question.

\subsubsection{1. The surface case, including nonorientable surfaces}
For an oriented closed Riemannian surface, Gauss--Bonnet gives
\[2\pi\chi(M)=\int_M K\,dA.\]
If $K\ge0$, the right side is nonnegative, proving the claim. If $M$ is nonorientable, its oriented double cover $\widetilde M$ inherits the lifted metric with the same local curvature sign. Euler characteristic is multiplicative under finite covers, since every cell of a finite CW structure has exactly the degree-many lifts. Hence $\chi(\widetilde M)=2\chi(M)$ and the oriented conclusion gives the desired result. No simply connected assumption enters.

\subsubsection{2. Products preserve both the curvature condition and the inequality}
For a product metric on $M\times N$, write two tangent vectors as $u=(u_M,u_N)$, $v=(v_M,v_N)$. The curvature numerator is
\[\langle R(u,v)v,u\rangle
=\langle R_M(u_M,v_M)v_M,u_M\rangle+
\langle R_N(u_N,v_N)v_N,u_N\rangle.\]
Each term is nonnegative when that factor has nonnegative sectional curvature; if its two projected vectors are dependent, that term is zero. Dividing by $|u\wedge v|^2>0$ proves nonnegative sectional curvature of every product plane, including mixed planes.

K\"unneth over $\mathbb Q$ gives $b_k(M\times N)=\sum_{i+j=k}b_i(M)b_j(N)$. Taking alternating sums of this finite identity yields
\[\chi(M\times N)=\chi(M)\chi(N).\]
Thus products of even-dimensional factors already satisfying the desired sign again satisfy it. An odd-dimensional closed manifold has Euler characteristic zero by Poincar\'e duality over an orientation cover (and division by its finite degree), so a product with an odd-dimensional closed factor has Euler characteristic zero as well. In particular, products of round spheres, flat tori and closed nonnegatively curved surfaces satisfy the target directly.

If a finite group of order $d$ acts freely and isometrically on such a product, the quotient inherits nonnegative sectional curvature because the projection is a local isometry. The same cellular covering count gives $\chi(\text{quotient})=\chi(\text{product})/d\ge0$. Freeness matters: a quotient with fixed points is an orbifold and need not be a manifold in the selected target.

\subsubsection{3. An explicit failed weakening of the hypothesis}
Let $\Sigma_g$ be a closed hyperbolic surface with genus $g>1$ and curvature $-1$, and put curvature $c>1$ on a round $S^2$ by rescaling its metric. The product four-manifold has scalar curvature
\[\mathrm{Scal}(S^2\times\Sigma_g)=2c-2>0,\]
since scalar curvatures add for products and a surface has scalar curvature twice its Gaussian curvature. Yet
\[\chi(S^2\times\Sigma_g)=2(2-2g)<0.\]
This is a concrete counterexample to a scalar-curvature version. It is not a counterexample to the question: tangent planes lying in the hyperbolic factor have sectional curvature $-1$. Therefore an attempted proof that replaces the full sectional-curvature operator by its scalar trace throws away indispensable information.

\subsubsection{4. Verification and remaining gap}
The denominator in the sectional-curvature formula is the area squared of the whole product plane, not the sum of factor areas, but only its positivity is needed for the sign argument. Euler-characteristic multiplicativity is independent of metric and follows from exact Betti-number algebra. These observations prove the stated closure families. In higher even dimension, the Gauss--Bonnet integrand is a polynomial in curvature components; nonnegative sectional curvature does not make its sign immediate from the two-dimensional argument. No pointwise sign theorem or alternative global argument is established here for a general manifold. That is the remaining obstruction.

\subsection{Sources}
\begin{itemize}
\item[S1]Manuel Amann. Bounds on Sectional Curvature and Interactions with Topology. 2020.
\url{https://link.springer.com/article/10.1365/s13291-020-00225-x}.
\textit{Formulation source.}
\end{itemize}
\end{document}
